\section{Markov chains}
\label{m:sec:chains}

A stochastic process is \emph{Markov} if its future depends on its past only
through its present. That single property is what makes every model in this thesis
tractable, and this section sets out the classes of Markov process the thesis uses.
Standard treatments are Norris \cite{Norris1997}, Karlin and Taylor
\cite{KarlinTaylor1975}, and, closest to the applications here, Allen
\cite{Allen2010}.

\subsection{Discrete-time Markov chains}
\label{m:sec:dtmc}

Let $\{X_n : n=0,1,2,\dots\}$ take values in a countable state space $\mathcal S$.
The process is a \emph{discrete-time Markov chain} if
\begin{equation}
\label{m:eq:dtmarkov}
\pr{X_{n+1}=j \mid X_n=i,\ X_{n-1},\dots,X_0} = \pr{X_{n+1}=j\mid X_n=i}
=: P_{ij},
\end{equation}
the right-hand side being independent of $n$ for a \emph{time-homogeneous} chain.
The matrix $P = (P_{ij})$ is stochastic: its entries are non-negative and each row
sums to one. Multi-step laws follow by decomposing a journey at an intermediate
state, which gives the Chapman--Kolmogorov identity
\begin{equation}
\label{m:eq:chapmankolmogorov}
\pr{X_{n+m}=j\mid X_0=i}
= \sum_{k\in\mathcal S}\pr{X_n=k\mid X_0=i}\,\pr{X_m=j\mid X_0=k},
\end{equation}
and hence $\pr{X_n=j\mid X_0=i}=(P^n)_{ij}$: composing steps is multiplying
matrices. In principle everything about the chain is therefore contained in $P$;
in practice the interest lies in extracting it, and the rest of this chapter is
about how. A state $i$ with $P_{ii}=1$ is \emph{absorbing}: a chain that reaches
it never leaves.

\subsubsection{Simple random walks}
\label{m:sec:randomwalk}

The simplest non-trivial chain moves one step at a time on the integers. Fix
$q\in(0,1)$ and let
\begin{equation}
\label{m:eq:rwrates}
P_{i,i+1} = q,
\qquad
P_{i,i-1} = 1-q,
\end{equation}
so that $X_n = X_0 + \sum_{k=1}^{n}\xi_k$ with the $\xi_k$ independent and equal to
$+1$ with probability $q$ and $-1$ otherwise. The walk is \emph{symmetric} when
$q=\tfrac12$. Its mean and variance are immediate from the independence of the
increments:
\begin{equation}
\label{m:eq:rwmoments}
\ex{X_n} = X_0 + n(2q-1),
\qquad
\var{X_n} = 4nq(1-q).
\end{equation}
The drift $2q-1$ decides the long-run behaviour, and the diffusive spread
$\sqrt{n}$ is what the mean cannot see --- the same tension, in its simplest form,
that motivates everything else in this chapter.

Here the whole law is available too, and it is worth writing down once, since it
is the distribution that \cref{m:eq:rwmoments} summarises. If $R$ of the $n$
increments are $+1$ then $X_n = X_0+2R-n$, and $R$ is binomial with parameters
$n$ and $q$, so
\begin{equation}
\label{m:eq:rwlaw}
\pr{X_n = X_0+2r-n} = \binom{n}{r}q^{\,r}(1-q)^{\,n-r},
\qquad r=0,1,\dots,n ,
\end{equation}
supported on those lattice points with the parity of $n$
(\cref{m:fig:rwlaw}). Almost no other chain in this chapter surrenders its
$n$-step law so cheaply, and the reason is that the increments of the walk are
independent of the state it has reached. Once the step distribution depends on
the current count --- as it does for every branching process below --- the
$n$-step law is no longer a convolution and must be reached by the generating
functions of \cref{m:sec:pgfmethod}.

\begin{figure}[ht]
\centering
\includegraphics[width=0.74\textwidth]{rw_transition.pdf}
\caption{The exact displacement law \cref{m:eq:rwlaw} at $n=40$, for
$q=0.4,\,0.5,\,0.6$; the dotted vertical lines mark the means $(2q-1)n$ of
\cref{m:eq:rwmoments}. Changing the drift slides the whole distribution along the
lattice and barely deforms it, because the width $\sqrt{4nq(1-q)}$ varies little
over this range of $q$. \Cref{m:fig:randomwalk} shows the same process as
trajectories; this is the distribution those trajectories are drawn from, and the
gap between the two pictures --- a smooth law, and paths that wander far from its
centre --- is the reason a mean is not an answer.}
\label{m:fig:rwlaw}
\end{figure}

The random walk is not a peripheral example here. Every continuous-time chain in
\cref{m:sec:ctmc} decomposes into a holding time and a jump, and the sequence of
jumps is itself a discrete-time chain --- the \emph{embedded jump chain}. For the
linear birth--death process of \cref{m:sec:bd}, whose rates out of state $i$ are
$\lambda i$ and $\mu i$, the state-dependence cancels in the ratio and the embedded
chain is exactly \cref{m:eq:rwrates} with
\begin{equation}
\label{m:eq:embedded}
q = \frac{\lambda}{\lambda+\mu} .
\end{equation}
Absorption at $0$ is then the walk's ruin problem, and hitting probabilities for
the continuous-time process can be computed on the walk alone; \cref{m:sec:firststep}
does exactly that. \Cref{m:eq:embedded} is also the bridge between the
discrete-generation and continuous-time models compared in
\cref{m:sec:condmeans}.

Two absorbing barriers are usually imposed. With absorbing states at $0$ and $M$,
writing $h_i$ for the probability of reaching $M$ before $0$ from state $i$, the
one-step decomposition $h_i = qh_{i+1}+(1-q)h_{i-1}$ with $h_0=0$, $h_M=1$ has the
classical gambler's-ruin solution
\begin{equation}
\label{m:eq:gamblersruin}
h_i =
\begin{cases}
\dfrac{1-\theta^{\,i}}{1-\theta^{\,M}},
& \theta := \dfrac{1-q}{q}\neq1,\\[12pt]
\dfrac{i}{M}, & q=\tfrac12 .
\end{cases}
\end{equation}
This is the smallest instance of the first-step principle of
\cref{m:sec:firststep}, and it is worth keeping in view: the entire method of
\cref{m:sec:methods} consists of conditioning on one step and solving the resulting
self-consistent equation, on state spaces where the equation is no longer a
two-term recursion. \Cref{m:fig:randomwalk} shows both halves of the picture.

\begin{figure}[ht]
\centering
\includegraphics[width=0.98\textwidth]{random_walk.pdf}
\caption{The simple random walk of \cref{m:eq:rwrates}.
(a) Seven realisations at $q=0.45$ (grey), against the drift line
$n(2q-1)$ and the envelope $\pm\sqrt{4nq(1-q)}$ of one standard deviation
(red), from \cref{m:eq:rwmoments}. The drift is linear in $n$ and the spread only
$O(\sqrt n)$, so the drift eventually dominates; over the first hundred steps it
does not, and individual paths spend long intervals on the wrong side of the mean.
(b) The gambler's-ruin probability \cref{m:eq:gamblersruin} of reaching $M=20$
before $0$, against the starting state $i$, for
$\theta=(1-q)/q \in\{0.7,\,0.85,\,1,\,1.2,\,1.5\}$. The symmetric case $\theta=1$
is the straight line $i/M$; $\theta<1$ bows the curve up and $\theta>1$ bows it
down, and the effect is strongly nonlinear in $\theta$. Only the ratio $\theta$
enters, which is why \cref{m:eq:bdruin} can read this figure off the embedded jump
chain of a birth--death process without reference to its rates.}
\label{m:fig:randomwalk}
\end{figure}

\subsubsection{Galton--Watson branching processes}
\label{m:sec:gw}

The Galton--Watson process is a discrete-time chain on the non-negative integers in
which every member of a population independently leaves offspring, or does not, at
each generational increment. It is the oldest model of stochastic branching and the
central object of this thesis; the standard references are Harris \cite{Harris1963}
and Athreya and Ney \cite{Athreya1972}.

Suppose each individual produces, independently of every other, a number of
offspring distributed as the \rv{} $L$, with $\pr{L=k}=p_k$ and mean
$\mu=\ex{L}$. Let $Z_n$ be the size of generation $n$, started from a single
progenitor, $Z_0=1$, so that
\begin{equation}
\label{m:eq:gwrecursion}
Z_{n+1} = \sum_{i=1}^{Z_n}L_i ,
\end{equation}
the $L_i$ being independent copies of $L$. Conditioning on the previous generation
and using the tower property gives $\ex{Z_{n+1}} = \mu\ex{Z_n}$, hence
\begin{equation}
\label{m:eq:gwmean}
\ex{Z_n} = \mu^{\,n} .
\end{equation}
When $\mu<1$ the expected population decays geometrically, and since $Z_n$ is
integer-valued Markov's inequality gives $\pr{Z_n\ge1}\le\mu^n\to0$: extinction is
certain. The three regimes $\mu<1$, $\mu=1$, $\mu>1$ are called
\emph{subcritical}, \emph{critical} and \emph{supercritical}.

Because $Z_{n+1}$ is a random sum of independent copies of $L$, the composition
rule for generating functions (\cref{m:sec:pgfmethod}) turns
\cref{m:eq:gwrecursion} into a functional iteration. Writing
$\phi(z)=\ex{z^L}$ for the offspring generating function and $\phi_n$ for that of
$Z_n$,
\begin{equation}
\label{m:eq:gwcomposition}
\phi_{n+1} = \phi_n\circ\phi ,
\end{equation}
so $\phi_n$ is the $n$-fold composition of $\phi$ with itself. That single fact is
the source of both the tractability and the difficulty of the model: everything
about $Z_n$ is encoded in an iterated map, and iterating a map is easy to do and
hard to solve in closed form.

\paragraph{The binary case.}
Throughout this thesis we specialise to the binary, or \emph{simple}, case, in
which an individual either divides into two or dies:
\begin{equation}
\label{m:eq:binarylaw}
p_2 = p \quad\text{(division)},
\qquad
p_0 = 1-p \quad\text{(death)},
\qquad
p_i = 0 \ \text{ otherwise},
\end{equation}
so that $\mu=2p$, $\ex{Z_n}=(2p)^n$, and the offspring generating function is
\begin{equation}
\label{m:eq:pgfbinary}
\phi(z) = pz^2+(1-p) .
\end{equation}
Criticality is at $p=\tfrac12$. Writing $u_n = \pr{Z_n=0}$ for the probability of
extinction by generation $n$ and setting $z=0$ in \cref{m:eq:gwcomposition} gives
$u_{n+1}=\phi(u_n)$, and it is more convenient to work with the survival
probability $S_n = 1-u_n = \pr{Z_n>0}$, which satisfies
\begin{equation}
\label{m:eq:SnIter}
S_{n+1} = 2pS_n - pS_n^2,
\qquad S_0 = 1 .
\end{equation}
\Cref{m:fig:gwvis} shows realisations either side of criticality.

\begin{figure}[ht]
\centering
\includegraphics[width=0.72\textwidth]{simpleGWvis.png}\\[6pt]
\includegraphics[width=0.72\textwidth]{subGWvis.png}
\caption{Realisations of the simple Galton--Watson process, drawn as genealogies
with generation increasing to the right. Top: supercritical, $p=0.58$, where the
number of lines grows without bound. Bottom: subcritical, $p=0.47$; the same
construction, but the tree terminates. Colour encodes generation number.}
\label{m:fig:gwvis}
\end{figure}

That is as far as the definition needs to go here. \Cref{m:eq:SnIter} is a quadratic
recursion, and almost everything of interest about the binary process turns out to
be a statement about its orbit: whether that orbit reaches zero
(\cref{m:sec:extinction}), how fast it does so (\cref{m:app:criticalgw}), and what
constant survives in front of the decay (\cref{m:sec:dtlimit}, and then \ChA{} at
length). The critical case, in which extinction is certain and yet the expected
lifetime is infinite, is developed in full in \cref{m:app:critical}.

\subsection{Continuous-time Markov chains}
\label{m:sec:ctmc}

Let $\{X_t : t\ge0\}$ take values in a countable state space $\mathcal S$, which
throughout will be a set of population counts or of tuples of counts. The process
is a \emph{continuous-time Markov chain} if, for all $s,t\ge0$,
\begin{equation}
\label{m:eq:markov}
\pr{X_{t+s}=j \mid X_t=i,\ \{X_u\}_{u<t}} = \pr{X_{t+s}=j\mid X_t=i},
\end{equation}
and \emph{time homogeneous} if the right-hand side depends on $s$ but not $t$, in
which case one writes $p_{ij}(s)$. Such a chain is specified by its
\emph{generator} $Q=(q_{ij})$, whose off-diagonal entries are transition rates,
\begin{equation}
\label{m:eq:generator}
q_{ij} = \lim_{\Delta t\to0}\frac{p_{ij}(\Delta t)}{\Delta t}\quad(j\neq i),
\qquad
q_{ii} = -\sum_{j\neq i}q_{ij} =: -q_i ,
\end{equation}
so that every row of $Q$ sums to zero. Equivalently, and this is the form used to
define models below,
\begin{equation}
\label{m:eq:infinitesimal}
\pr{X_{t+\Delta t}=j\mid X_t=i} = q_{ij}\,\Delta t + o(\Delta t)
\qquad(j\neq i),
\end{equation}
with complementary probability $1-q_i\Delta t+o(\Delta t)$ of no transition at all.
A state with $q_i=0$ is absorbing.

\subsubsection{Exponential clocks}
\label{m:sec:exp}

Everything in continuous time rests on one distribution, and it is worth seeing
why it is forced rather than chosen. A \rv{} $T$ is exponentially distributed with
rate $\theta>0$, written $T\sim\mathrm{Exp}(\theta)$, if
\begin{equation}
\label{m:eq:expdef}
f_T(t) = \theta\ee^{-\theta t},
\qquad
F_T(t) = 1-\ee^{-\theta t}
\qquad(t\ge0),
\end{equation}
with mean $1/\theta$ and variance $1/\theta^2$. The property that matters is
memorylessness: for $t_1,t_2>0$,
\begin{equation}
\label{m:eq:memoryless}
\pr{T>t_1+t_2\mid T>t_1}
= \frac{\ee^{-\theta(t_1+t_2)}}{\ee^{-\theta t_1}}
= \ee^{-\theta t_2}
= \pr{T>t_2},
\end{equation}
so a clock that has already run for time $t_1$ is statistically indistinguishable
from a fresh one. The exponential law is the only continuous distribution with this
property, which is exactly why a process assembled from constant-rate events
satisfies \cref{m:eq:markov}: with any other waiting-time law the elapsed time would
be informative and the process would not be Markov.

When several events compete, two facts organise everything that follows.

\begin{proposition}[Competing clocks]
\label{m:prop:race}
Let $T_1,\dots,T_m$ be independent with $T_i\sim\mathrm{Exp}(\theta_i)$, let
$T=\min_iT_i$, and write $\Theta=\sum_i\theta_i$. Then
\begin{equation}
\label{m:eq:race}
T\sim\mathrm{Exp}(\Theta),
\qquad
\pr{T_k=T} = \frac{\theta_k}{\Theta},
\end{equation}
and the identity of the winning clock is independent of $T$.
\end{proposition}

\begin{proof}
For the first claim,
$\pr{T>t} = \prod_i\pr{T_i>t} = \prod_i\ee^{-\theta_it} = \ee^{-\Theta t}$.
For the second, condition on the value of $T_k$ and use independence:
\begin{equation*}
\pr{T_k=T}
= \int_0^\infty\theta_k\ee^{-\theta_ku}\prod_{i\neq k}\ee^{-\theta_iu}\,\dd u
= \theta_k\int_0^\infty\ee^{-\Theta u}\,\dd u
= \frac{\theta_k}{\Theta} .
\end{equation*}
Repeating the calculation with the integral truncated at $t$ gives
$\pr{T_k=T,\ T>t} = (\theta_k/\Theta)\ee^{-\Theta t} = \pr{T_k=T}\pr{T>t}$, which is
the asserted independence.
\end{proof}

\Cref{m:prop:race} is what connects \cref{m:eq:infinitesimal} to a trajectory: the
chain leaves state $i$ after a holding time $\mathrm{Exp}(q_i)$ and then jumps to
$j$ with probability $q_{ij}/q_i$, independently of how long it waited. It also
says that a rate becomes a probability as soon as one asks \emph{which} of several
things happened rather than when, which is the identification \cref{m:eq:embedded}
of the embedded jump chain. And it is the whole content of the direct simulation
method.

\begin{remark}[Direct simulation]
\label{m:rem:gillespie}
\Cref{m:prop:race} is an algorithm as well as a theorem. Given the current state,
compute the rate $\theta_i$ of every possible event and their sum $\Theta$; draw
$T\sim\mathrm{Exp}(\Theta)$ by inverting \cref{m:eq:expdef} on a uniform variate;
draw the event $k$ with probability $\theta_k/\Theta$; advance the clock by $T$,
apply event $k$, and repeat. Because the two draws are independent and the
exponential law is memoryless, the resulting trajectory has exactly the law of the
process. This is Gillespie's direct method \cite{Gillespie1977}; it is used in this
thesis to check analytical results, never to replace them.
\end{remark}

\subsubsection{The Poisson process}
\label{m:sec:poisson}

The simplest continuous-time chain counts events that arrive at a constant rate and
never depart. Let $N_t$ take values in $\{0,1,2,\dots\}$ with the single non-zero
rate
\begin{equation}
\label{m:eq:poissonrates}
q_{n,n+1} = \alpha
\qquad(n\ge0),
\end{equation}
so $q_n=\alpha$ for every $n$. By \cref{m:prop:race} the holding times are
independent $\mathrm{Exp}(\alpha)$ variables, and $N_t$ counts how many have
elapsed by time $t$: the process is the renewal count of a sequence of independent
exponential waits. Its law follows from the master equation
\cref{m:eq:master} below, or directly by induction on $n$:
\begin{equation}
\label{m:eq:poissonlaw}
\pr{N_t = n} = \ee^{-\alpha t}\frac{(\alpha t)^n}{n!},
\qquad
\ex{N_t} = \var{N_t} = \alpha t .
\end{equation}
The equality of mean and variance is characteristic, and is the baseline against
which the over-dispersion of branching processes should be read.
\Cref{m:fig:poisson} shows a realisation and the law it is drawn from.

\begin{figure}[ht]
\centering
\includegraphics[width=0.98\textwidth]{poisson_process.pdf}
\caption{The Poisson process of \cref{m:eq:poissonrates}, at rate $\alpha=1.5$.
(a) One realisation of the counting process $N_t$, with the independent
$\mathrm{Exp}(\alpha)$ holding times of \cref{m:prop:race} marked in alternating
shades along the axis, and the mean $\alpha t$ dashed. The path crosses and
recrosses its own mean; nothing prevents a long wait followed by three arrivals in
quick succession, because each holding time is drawn afresh.
(b) The law \cref{m:eq:poissonlaw} at $t=1,\,3,\,6$. Mean and variance are both
$\alpha t$, so the distribution broadens as it advances and its relative width
$1/\sqrt{\alpha t}$ shrinks.}
\label{m:fig:poisson}
\end{figure}

Two properties are used later. Counts over disjoint time intervals are
independent, and the count over an interval depends only on its length --- the
process has stationary independent increments, which is the continuous-time
statement of memorylessness. And if each arrival is independently retained with
probability $\varpi$, the retained arrivals form a Poisson process of rate
$\varpi\alpha$: \emph{thinning} preserves the class. Thinning is what makes the
absorption models of \cref{m:sec:moc} tractable, where a cohort of exterior
particles is taken up one by one at a constant per-particle rate.

The Poisson process is also the $\mu=0$, state-independent corner of the family
treated next. Making the arrival rate proportional to the current count, $\alpha n$
in place of $\alpha$, gives the pure-birth (Yule) process; adding a departure rate
gives the birth--death process.

\subsubsection{Birth--death processes}
\label{m:sec:bd}

Let $X_t$ count individuals, each of which independently gives birth at rate
$\lambda$ and dies at rate $\mu$. The non-zero rates out of state $n$ are
\begin{equation}
\label{m:eq:bdrates}
q_{n,n+1} = \lambda n,
\qquad
q_{n,n-1} = \mu n ,
\end{equation}
so $q_n = (\lambda+\mu)n$ and $n=0$ is absorbing. This is the \emph{linear}
birth--death process, the continuous-time counterpart of the Galton--Watson process
of \cref{m:sec:gw}, and it is the running example of \cref{m:sec:ctmethods}: nearly
every method in this thesis is first demonstrated on it. Its embedded jump chain is
the random walk \cref{m:eq:rwrates} with $q=\lambda/(\lambda+\mu)$, so questions
about \emph{which} states are visited reduce to \cref{m:sec:randomwalk} while
questions about \emph{when} require the rates themselves.

\Cref{m:fig:bdpaths} shows the three regimes it can occupy.

\begin{figure}[ht]
\centering
\includegraphics[width=0.98\textwidth]{birth_death_paths.pdf}
\caption{Nine realisations of the linear birth--death process \cref{m:eq:bdrates}
from $X_0=5$, with $\mu=1$ throughout and the mean $N\ee^{(\lambda-\mu)t}$ of
\cref{m:eq:bdmean} in red; dots on the axis mark trajectories that reached the
absorbing state. The middle panel is the one to dwell on. At criticality the mean
is exactly constant, and yet several realisations are already extinct while others
have grown fourfold --- the mean is the same number in both cases, and describes
neither. That gap between a mean and a distribution is what the conditioning of
\cref{m:sec:condmeans} exists to close. Note also that extinction is not confined
to the subcritical panel: a supercritical process still dies out with probability
$(\mu/\lambda)^{N}$, visible here as the paths that fail early on the right.}
\label{m:fig:bdpaths}
\end{figure}

The general birth--death process allows state-dependent rates $\lambda_n$ and
$\mu_n$; the linear case $\lambda_n=\lambda n$, $\mu_n=\mu n$ is the one in which
individuals act independently, and hence the one for which generating functions
close. Everything derived in \cref{m:sec:ctmethods} --- mean, variance, extinction
probability, conditional mean --- uses that independence.

\subsubsection{Birth--death--catastrophe processes}
\label{m:sec:bdc}

The processes at the centre of this thesis add a third channel. Alongside the
gradual loss of individuals to death, a compartment may suffer a
\emph{catastrophe}: an abrupt event removing the entire population at once.

\begin{definition}[Birth--death--catastrophe process]
\label{m:def:bdc}
Let $X_t$ count individuals in a compartment. Each individual independently gives
birth at rate $\lambda$ and dies at rate $\mu$, and the compartment additionally
undergoes a catastrophe --- a rupture that removes the entire population at once
--- at rate $\rho$ per individual, so that the total catastrophe rate in state $n$
is $\rho n$. Writing $\dagger$ for the post-catastrophe state, the non-zero
transition rates out of $n\ge1$ are
\begin{equation}
\label{m:eq:bdcrates}
q_{n,n+1} = \lambda n,
\qquad
q_{n,n-1} = \mu n,
\qquad
q_{n,\dagger} = \rho n ,
\end{equation}
and both $n=0$ and $\dagger$ are absorbing.
\end{definition}

Two features distinguish \cref{m:def:bdc} from the birth--death process, and both
matter later. The process has \emph{two} absorbing states, reached by mechanisms
that are distinguishable and not interchangeable: extinction empties the
compartment gradually, catastrophe empties it in one event and releases its
contents. And the catastrophe rate is proportional to the population, so the
compartment becomes more likely to rupture the fuller it is --- the hazard is
state-dependent, which is a genuine complication rather than a parameter choice, as
\cref{m:rem:notmeanfield} shows. \Cref{m:fig:ratediagram} sets the two processes
side by side.

\begin{figure}[ht]
\centering
\begin{tikzpicture}[
  >={Stealth[length=1.9mm]},
  st/.style ={circle, draw, thick, minimum size=7.5mm, inner sep=0pt, font=\small},
  ab/.style ={circle, draw, thick, double, double distance=0.7pt,
              minimum size=7.5mm, inner sep=0pt, font=\small},
  lb/.style ={font=\footnotesize, inner sep=1.6pt},
  rowlab/.style={font=\small\itshape, anchor=west},
  every path/.style={thick}
]
% ---------------- linear birth--death -------------------------------------
\node[rowlab] at (-1.9,1.15) {birth--death};
\node[ab] (a0) at (0,0)   {$0$};
\node[st] (a1) at (2.1,0) {$1$};
\node[st] (a2) at (4.2,0) {$2$};
\node[st] (a3) at (6.3,0) {$3$};
\node     (a4) at (7.9,0) {$\cdots$};
\draw[->] (a1) to[bend left=38]  node[above,lb]{$\lambda$}  (a2);
\draw[->] (a2) to[bend left=38]  node[above,lb]{$2\lambda$} (a3);
\draw[->] (a3) to[bend left=38]  node[above,lb]{$3\lambda$} (a4);
\draw[->] (a1) to[bend left=38]  node[below,lb]{$\mu$}      (a0);
\draw[->] (a2) to[bend left=38]  node[below,lb]{$2\mu$}     (a1);
\draw[->] (a3) to[bend left=38]  node[below,lb]{$3\mu$}     (a2);

% ---------------- birth--death--catastrophe -------------------------------
\begin{scope}[yshift=-3.5cm]
\node[rowlab] at (-1.9,1.15) {birth--death--catastrophe};
\node[ab] (b0) at (0,0)   {$0$};
\node[st] (b1) at (2.1,0) {$1$};
\node[st] (b2) at (4.2,0) {$2$};
\node[st] (b3) at (6.3,0) {$3$};
\node     (b4) at (7.9,0) {$\cdots$};
\draw[->] (b1) to[bend left=38]  node[above,lb]{$\lambda$}  (b2);
\draw[->] (b2) to[bend left=38]  node[above,lb]{$2\lambda$} (b3);
\draw[->] (b3) to[bend left=38]  node[above,lb]{$3\lambda$} (b4);
\draw[->] (b1) to[bend left=38]  node[below,lb]{$\mu$}      (b0);
\draw[->] (b2) to[bend left=38]  node[below,lb]{$2\mu$}     (b1);
\draw[->] (b3) to[bend left=38]  node[below,lb]{$3\mu$}     (b2);
\node[ab] (bd) at (4.2,-2.35) {$\dagger$};
\draw[->,densely dashed] (b1) to[bend right=12]
      node[pos=0.68,left=1.5pt,lb]{$\rho$} (bd);
\draw[->,densely dashed] (b2) --
      node[pos=0.68,right=1.5pt,lb]{$2\rho$} (bd);
\draw[->,densely dashed] (b3) to[bend left=12]
      node[pos=0.68,right=1.5pt,lb]{$3\rho$} (bd);
\end{scope}
\end{tikzpicture}
\caption{Transition-rate diagrams for the linear birth--death process
\cref{m:eq:bdrates} and for the birth--death--catastrophe process of
\cref{m:def:bdc}. Double circles mark absorbing states. Every rate is proportional
to the current count, because individuals act independently. The two differ in one
structural respect: the birth--death chain can only be absorbed by stepping down
through state $1$, whereas catastrophe (dashed) removes the population from
\emph{any} state in a single transition, and does so faster the fuller the
compartment is. That is why survival mass drains from state $1$ alone in the first
process but from every positive state in the second, as
\cref{m:eq:bdcsurvivalloss} records, and it is what makes $\dagger$ and $0$
genuinely different endings rather than two labels for extinction.}
\label{m:fig:ratediagram}
\end{figure}

It is useful to see the whole family at once. Let the positive state $i$ carry
birth rate $\lambda_i$, ordinary death rate $\mu_i$ and catastrophe rate
$\kappa_i$, a catastrophe sending the process directly to $0$. On bounded test
functions the generator is
\begin{equation}
\label{m:eq:bdcgenerator}
(\mathcal Lf)(i)
= \lambda_i\{f(i+1)-f(i)\} + \mu_i\{f(i-1)-f(i)\} + \kappa_i\{f(0)-f(i)\},
\qquad i\ge1,
\end{equation}
with $(\mathcal Lf)(0)=0$. Restricting to the positive states gives the
\emph{killed} subgenerator $K$, with $K_{i,i+1}=\lambda_i$, $K_{i,i-1}=\mu_i$ for
$i\ge2$, and
\begin{equation}
\label{m:eq:killedsubgen}
K_{i,i} = -(\lambda_i+\mu_i+\kappa_i) .
\end{equation}
The row mass missing from $K$ is the instantaneous absorption rate: $\kappa_i$ for
$i\ge2$, and $\mu_1+\kappa_1$ for $i=1$. Writing $T_0=\inf\{t\ge0:X_t=0\}$ and
summing the positive-state forward equations,
\begin{equation}
\label{m:eq:bdcsurvivalloss}
\ddt{}\pr{T_0>t}
= -\mu_1\pr{X_t=1} - \sum_{i\ge1}\kappa_i\pr{X_t=i},
\end{equation}
all probabilities on the right being understood before absorption. Ordinary death
drains survival mass only from state $1$; catastrophe drains it from every positive
state, and that asymmetry is the whole difference between the two channels. The
consequences for conditioning are taken up in \cref{m:sec:condmeans}, and the
process itself is developed \fwd{bdc}{in the later chapters on
birth--death--catastrophe processes}, with the two-type extension treated
\fwd{multitype}{in the chapter on multi-type processes} and the release of a
ruptured compartment's contents into a shared medium
\fwd{rupture}{in the chapter on compartment rupture}.

\begin{remark}[A state-dependent catastrophe is not a mean-field hazard]
\label{m:rem:notmeanfield}
It is tempting to replace a state-dependent catastrophe rate by a deterministic
hazard evaluated at the mean population, and it does not work. The point is
clearest in discrete time. Let $\kappa$ be the probability that a given individual
initiates catastrophe at a step, let $\widetilde Z_i$ be the branching population
with catastrophe switched off, and let $H_n=\sum_{i=0}^{n-1}\widetilde Z_i$ be the
cumulative exposure. Conditioning on the entire catastrophe-free genealogy, every
individual present independently declines to trigger, so
\begin{equation}
\label{m:eq:norupture}
\pr{\text{no catastrophe before }n} = \ex{(1-\kappa)^{H_n}} .
\end{equation}
This identity is exact. Replacing $H_n$ by $\ex{H_n}$ is not: the map
$h\mapsto(1-\kappa)^h$ is convex, so Jensen's inequality gives
$\ex{(1-\kappa)^{H_n}}\ge(1-\kappa)^{\ex{H_n}}$, and the substitution
systematically understates the chance of getting through. The failure is
structural, not a matter of accuracy --- the expectation in \cref{m:eq:norupture} is
a nonlinear functional of the whole path, and no deterministic hazard closes it.
Two further cautions belong with it. Non-catastrophe is not the same event as
population survival: the probability of being both unruptured and non-zero at
generation $n$ is $\ex{\ind\{\widetilde Z_n>0\}(1-\kappa)^{H_n}}$, for which
Jensen supplies no bound at all. And once a trigger occurs every surviving lineage
is stopped together, so the killed lineages are coupled and cannot be treated as
independent after the fact --- although their family trees and trigger variables do
factorise \emph{before} any trigger, which is what makes \cref{m:eq:norupture}
correct. Closing the model properly therefore requires a marked or multivariate
generating function, or the killed-generator viewpoint of \cref{m:eq:killedsubgen};
the elementary calculation this remark displaces, and the size of the gap it leaves,
are recorded in \cref{m:app:discreterupture}.
\end{remark}

\subsection{Time-inhomogeneous processes}
\label{m:sec:inhomogeneous}

All the chains above are time homogeneous: the rates in \cref{m:eq:infinitesimal}
depend on the state but not on the clock. Dropping that assumption gives a
\emph{time-inhomogeneous} chain, with generator $Q(t)$ and transition
probabilities depending on both endpoints,
\begin{equation}
\label{m:eq:inhomogtrans}
p_{ij}(s,t) = \pr{X_t=j\mid X_s=i}
\qquad(s\le t),
\end{equation}
in place of the single-argument $p_{ij}(t-s)$. The Chapman--Kolmogorov identity
survives in the form $P(s,t) = P(s,u)P(u,t)$ for $s\le u\le t$, and differentiating
it gives forward and backward equations
\begin{equation}
\label{m:eq:inhomogkolmogorov}
\pdd{}{t}P(s,t) = P(s,t)\,Q(t),
\qquad
\pdd{}{s}P(s,t) = -Q(s)\,P(s,t),
\end{equation}
but the solution is no longer the matrix exponential $\ee^{(t-s)Q}$ unless the
generators at different times commute. In general one is left with a
time-ordered product, and the methods of \cref{m:sec:ctmethods} apply only after
that complication has been handled.

Much survives nonetheless. The competing-clocks decomposition of
\cref{m:prop:race} still holds instantaneously, so direct simulation still applies
with a time-dependent total rate --- at the cost of thinning or of integrating the
rate to find the next event time. And a linear birth--death process with
time-varying rates $\lambda(t)$, $\mu(t)$ still has a generating function
satisfying a first-order partial differential equation, so the method of
characteristics of \cref{m:sec:moc} remains available; the characteristics simply
acquire time-dependent coefficients.

\subsubsection{The time-inhomogeneous linear birth--death process}
\label{m:sec:tibd}

The workhorse of the inhomogeneous setting is the process of \cref{m:sec:bd} with
rates that are prescribed functions of the clock: each individual gives birth at
rate $\lambda(t)$ and dies at rate $\mu(t)$, both taken locally integrable. The
master equation is \cref{m:eq:bdmaster} with time-dependent coefficients, and the
generating function accordingly satisfies
\begin{equation}
\label{m:eq:tibdpde}
\pddt{G} = \bigl(\lambda(t)z^2-(\lambda(t)+\mu(t))z+\mu(t)\bigr)\pdd{G}{z},
\qquad
G(z,0)=z^N ,
\end{equation}
which is \cref{m:eq:bdpde} with its quadratic coefficient set in motion. The
equation is still first order, so \cref{m:sec:moc} still applies; what has been
lost is that the coefficient no longer factorises once and for all, and with it
the closed solution \cref{m:eq:bdsolution}.

Two quantities survive intact whatever the rates do. The mean obeys the same
closed equation as before, $\ddt{}\ex{X_t}=(\lambda(t)-\mu(t))\ex{X_t}$, because
the rates are still linear in the count, so
\begin{equation}
\label{m:eq:meanodeinhom}
\ex{X_t} = N\exp\lrb{\int_0^t\bigl(\lambda(s)-\mu(s)\bigr)\dd s} ,
\end{equation}
and only the integrated net growth enters, not how it was distributed over the
interval. And the extinction probability $u(t)=\pr{X_t=0}$ from a single founder
satisfies the Riccati equation
\begin{equation}
\label{m:eq:riccatiext}
\dda{u}{t} = \mu(t)-\bigl(\lambda(t)+\mu(t)\bigr)u+\lambda(t)u^2,
\qquad u(0)=0 ,
\end{equation}
which is the backward equation of \cref{m:sec:backward} with time-dependent
coefficients: over the first small interval either the founder dies, contributing
$\mu(t)\Delta t$, or it divides, contributing $\lambda(t)\Delta t\,u(t)^2$ since
both daughter lineages must then die out in turn, or nothing happens. This is a
single scalar equation, integrated numerically without difficulty when it does
not solve in closed form, and the full generating function retains the
fractional-linear shape of \cref{m:eq:bdsolution} with its coefficients replaced
by solutions of \cref{m:eq:riccatiext}.

The same Riccati structure recurs twice more below, and the three occurrences are
worth keeping apart. It appears as the continuum approximation
\cref{m:eq:riccati} to the discrete survival recursion, and as the characteristic
equation \cref{m:eq:ch23} of the absorption--birth--death model. Only the middle
one is an approximation; \cref{m:eq:riccatiext} and \cref{m:eq:ch23} are exact
consequences of their models, and the difference is exactly what
\cref{m:app:riccati} turns on.

\subsubsection{The logistic speciation model}
\label{m:sec:logisticspeciation}

The standard application of that framework is a population whose per-capita birth
rate falls as the population fills the space available to it. In the macroevolutionary
reading of that idea the counted individuals are species and the births are
speciation events \cite{Allen2010}: let $X_t$ be the standing diversity, let each
species go extinct at constant rate $\mu$, and let it speciate at the
diversity-dependent rate
\begin{equation}
\label{m:eq:speciationrate}
\lambda(n) = \lambda_0\lrb{1-\frac{n}{K}},
\qquad \lambda_0>\mu ,
\end{equation}
when the standing diversity is $n$, with $K$ the niche capacity.

\Cref{m:eq:speciationrate} is not yet of the form treated above, and the
distinction is the whole point of the subsection: the rate depends on the
\emph{state}, not on the clock. The chain it defines is nonlinear, its moment
equations do not close, and none of the machinery of \cref{m:sec:ctmethods}
reaches it directly. The tractable surrogate replaces the random diversity inside
\cref{m:eq:speciationrate} by a deterministic mean-field trajectory $\bar N(t)$,
which turns the rate into a prescribed function of time and the chain into a
time-inhomogeneous \emph{linear} process, to which \cref{m:sec:tibd} does apply.
That replacement is a modelling step and not an identity. It is the same
substitution of a mean for a random census that \cref{m:rem:notmeanfield} shows to
fail outright for a state-dependent catastrophe, and what separates the two cases is
only what is being asked of it: here it defines a tractable comparison process, there
it was asked to return an exact survival probability and did not. Everything that
follows is accordingly a statement about the surrogate, not about the chain
\cref{m:eq:speciationrate} defines.

The mean-field diversity solves the logistic equation
\begin{equation}
\label{m:eq:logisticode}
\dda{\bar N}{t} = \bigl(\lambda(\bar N)-\mu\bigr)\bar N
= \gamma\,\bar N\lrb{1-\frac{\bar N}{N_\ast}},
\qquad
\gamma = \lambda_0-\mu,
\quad
N_\ast = K\lrb{1-\frac{\mu}{\lambda_0}},
\end{equation}
whose equilibrium $N_\ast$ is the diversity at which speciation balances
extinction. From $\bar N(0)=N_0$,
\begin{equation}
\label{m:eq:logisticsol}
\bar N(t) = \frac{N_\ast}{1+B\ee^{-\gamma t}},
\qquad
B = \frac{N_\ast}{N_0}-1 ,
\end{equation}
and the surrogate speciation rate $\lambda(t)=\lambda_0\bigl(1-\bar N(t)/K\bigr)$
is
\begin{equation}
\label{m:eq:lambdat}
\lambda(t) = \mu + \frac{\gamma B\ee^{-\gamma t}}{1+B\ee^{-\gamma t}},
\qquad
\lambda(0) = \lambda_0\lrb{1-\frac{N_0}{K}} ,
\end{equation}
relaxing monotonically from $\lambda(0)$ down to $\mu$. The integrated net growth
that \cref{m:eq:meanodeinhom} calls for is then available in closed form,
\begin{equation}
\label{m:eq:lambdaint}
\int_0^t\bigl(\lambda(s)-\mu\bigr)\dd s
= \ln\lrb{\frac{1+B}{1+B\ee^{-\gamma t}}} ,
\end{equation}
so that the mean diversity of the surrogate process is
\begin{equation}
\label{m:eq:logspecmean}
\ex{X_t} = N_0\,\frac{1+B}{1+B\ee^{-\gamma t}}
\;\xrightarrow[t\to\infty]{}\;
N_0(1+B) = N_\ast ,
\end{equation}
recovering the deterministic equilibrium, as a linear surrogate built on that
equilibrium must. Its extinction probability follows by integrating
\cref{m:eq:riccatiext} with the rate \cref{m:eq:lambdat}, and no longer in closed
form. \Cref{m:fig:logspecmean} plots \cref{m:eq:logspecmean} for three founding
diversities.

\begin{figure}[ht]
\centering
\includegraphics[width=0.78\textwidth]{logspec_mean.pdf}
\caption{Mean diversity \cref{m:eq:logspecmean} of the logistic speciation
surrogate for founding diversities $N_0=1,3,6$, with $\mu/\lambda_0=0.1$ and
$K=10$, so that $N_\ast=9$; the abscissa is the rescaled time $\gamma t$. Every
founding diversity relaxes monotonically to $N_\ast$, and the mean of the
surrogate coincides with the deterministic solution \cref{m:eq:logisticsol}
exactly, because the surrogate is linear and its mean therefore closes. The
smallest cohort spends a long early interval in near-exponential growth before the
diversity dependence binds --- and it is over precisely that interval that the
surrogate is least trustworthy, since it is there that a single realisation is
least like its mean.}
\label{m:fig:logspecmean}
\end{figure}

Two limits are worth reading off \cref{m:eq:lambdat}. At early times
$\lambda(t)\approx\lambda(0)>\mu$ and the surrogate is indistinguishable from a
homogeneous supercritical process; at late times $\lambda(t)\downarrow\mu$ and it
approaches criticality from above. The model therefore drifts, over its own
lifetime, into the regime whose discrete-generation counterpart causes the
difficulties documented in \cref{m:app:criticalgw} --- a fixed carrying capacity
manufactures a near-critical process out of a comfortably supercritical one, and
does so without any parameter being changed by hand.

\subsection{Coupled ODE--CTMC systems}
\label{m:sec:coupled}

The final class leaves the purely stochastic setting. In a \emph{coupled
ODE--CTMC system} a continuous, deterministic state $\mathbf y(t)\in\mathbb R^d$
evolves alongside a discrete, stochastic state $X_t\in\mathcal S$, and the two
influence one another. Three couplings are possible, and they are genuinely
different.

\begin{enumerate}[label=(\roman*),leftmargin=*]
\item \emph{The differential equation depends on the chain.} The vector field is
      indexed by the discrete state,
      \begin{equation}
      \label{m:eq:coupledone}
      \dda{\mathbf y}{t} = F\bigl(\mathbf y(t),\,X_t\bigr),
      \end{equation}
      with $X_t$ an autonomous Markov chain. Between jumps the trajectory follows a
      fixed flow; at each jump the flow switches. The pair
      $(\mathbf y(t),X_t)$ is a Markov process on $\mathbb R^d\times\mathcal S$ even
      though neither component is one alone, and it is a piecewise-deterministic
      Markov process in the standard classification \cite{Davis1993}.
\item \emph{The chain depends on the differential equation.} The transition rates
      are functions of the continuous state,
      \begin{equation}
      \label{m:eq:coupledtwo}
      q_{ij} = q_{ij}\bigl(\mathbf y(t)\bigr),
      \end{equation}
      with $\mathbf y$ following a prescribed trajectory. This is a
      time-inhomogeneous chain in the sense of \cref{m:sec:inhomogeneous}, with the
      inhomogeneity generated by the flow rather than imposed directly, and the
      remarks there apply.
\item \emph{Both directions at once.} The vector field depends on the discrete
      state and the rates depend on the continuous one:
      \begin{equation}
      \label{m:eq:coupledthree}
      \dda{\mathbf y}{t} = F\bigl(\mathbf y(t),X_t\bigr),
      \qquad
      q_{ij} = q_{ij}\bigl(\mathbf y(t)\bigr) .
      \end{equation}
      Neither component can now be solved first and substituted into the other, and
      the joint process must be treated as a whole.
\end{enumerate}

Case (i) alone is already enough to break the methods of \cref{m:sec:ctmethods} as
stated, because the state space is no longer countable and the generator acquires a
transport term: for suitable test functions $f(\mathbf y,i)$,
\begin{equation}
\label{m:eq:coupledgenerator}
(\mathcal Lf)(\mathbf y,i)
= F(\mathbf y,i)\cdot\nabla_{\mathbf y}f(\mathbf y,i)
+ \sum_{j\neq i}q_{ij}(\mathbf y)\bigl\{f(\mathbf y,j)-f(\mathbf y,i)\bigr\} .
\end{equation}
The first term is the deterministic drift along the flow and the second is the
familiar jump part; \cref{m:eq:coupledgenerator} covers all three cases above,
with the $\mathbf y$-dependence of the rates dropped in case (i) and the
$i$-dependence of $F$ dropped in case (ii). Written this way the coupled system is
a recognisable object, and the corresponding forward equation is a system of
transport equations indexed by $\mathcal S$ rather than a system of ordinary
differential equations.

\Cref{m:fig:coupled} shows what the two-way case \cref{m:eq:coupledthree} looks like
in a trajectory. The system drawn there is generic --- it is chosen only to make the
mechanism visible, and is not a model used elsewhere in this thesis --- but it
displays the feature that distinguishes coupling (iii) from the two one-way cases.
Because the switching rate rises with $y$, the intervals spent in the second
discrete state are not scattered uniformly in time: they cluster where the
continuous variable has climbed, and each such interval then pulls $y$ back down.
Neither component can be solved first and substituted into the other, because each
is continually re-shaping the other's behaviour.

\begin{figure}[ht]
\centering
\includegraphics[width=0.98\textwidth]{coupled_ode_ctmc.pdf}
\caption{A trajectory of a coupled ODE--CTMC system of the two-way form
\cref{m:eq:coupledthree}. The continuous variable follows
$\dd y/\dd t = a(X_t)-by$, so it relaxes towards $a_1/b$ while $X_t=1$ and towards
$a_2/b$ while $X_t=2$ (dashed lines); the chain switches $1\to2$ at rate
$\kappa y$, which depends on the continuous state, and back at a constant rate.
Upper panel: $y(t)$, with the intervals spent in state $2$ shaded. Lower panel:
the discrete component $X_t$. The visits to state $2$ concentrate where $y$ is
large, which is the $y$-dependence of the rate made visible; the resulting descents
then remove the very condition that made them likely. The joint process
$(y(t),X_t)$ is Markov although neither component is, and its generator is
\cref{m:eq:coupledgenerator}. The parameters are illustrative only.}
\label{m:fig:coupled}
\end{figure}

\subsubsection{Simulating a coupled system}
\label{m:sec:coupledsim}

The competing-clocks construction of \cref{m:prop:race} extends to
\cref{m:eq:coupledthree} with one qualification: the rates are no longer constant
between events, because $\mathbf y$ moves. Given the current state
$(\mathbf y,i)$, integrate the flow with the discrete component held at $i$ to
obtain $\mathbf y_s$ for $s\ge t$; the total jump rate
$q_i(\mathbf y_s)=\sum_{j\neq i}q_{ij}(\mathbf y_s)$ is then a known function of
$s$, and the next jump time $T$ has survival function
\begin{equation}
\label{m:eq:coupledholding}
\pr{T>s} = \exp\lrb{-\int_t^{s}q_i(\mathbf y_u)\,\dd u} ,
\end{equation}
exactly as in \cref{m:sec:inhomogeneous}; the destination is then drawn with
probability $q_{ij}(\mathbf y_T)/q_i(\mathbf y_T)$, the state is updated, and the
procedure repeats. Where $q_i$ admits a constant majorant the integral can be
inverted by rejection, and otherwise by quadrature. Analytically the position is
less comfortable. If the coupling is weak one closes the system by replacing
$\mathbf y$ by its mean-field trajectory --- which returns a time-inhomogeneous
chain in the sense of \cref{m:sec:inhomogeneous}, and the logistic surrogate of
\cref{m:sec:logisticspeciation} is exactly that closure --- or by replacing $X_t$
by its mean, which returns a closed differential equation. When the coupling is
strong neither closure is exact, and the working objects become joint
expectations satisfying hybrid moment equations that mix the jump generator with
the drift.

\subsubsection{A schematic instance: rupture into a shared medium}
\label{m:sec:coupledexample}

It is worth seeing the framework carry one concrete mechanism, and the natural
one here is suggested by \cref{m:def:bdc}. Let $X_t$ be the contents of a
compartment, evolving as the birth--death--catastrophe process of that
definition, and let $y(t)\ge0$ be the concentration of material in a medium the
compartment sits in. Between catastrophes the medium clears at a constant
per-unit rate $\delta$,
\begin{equation}
\label{m:eq:mediumdecay}
\dda{y}{t} = -\delta y ,
\end{equation}
and at a catastrophe time $\tau$ the entire contents of the compartment are
released into the medium at once,
\begin{equation}
\label{m:eq:mediumjump}
y(\tau+) = y(\tau-) + c\,X_{\tau-} ,
\end{equation}
with $c>0$ the material carried per individual. That is the first arrow of the
coupling: the discrete component drives the flow. The second is equally present,
since the catastrophe rate $\rho$, and any rate at which a surviving compartment
takes material back up, may themselves be functions of the current concentration,
so that a fuller medium hastens the next rupture. Both dependences are then
active and the system is of the two-way form \cref{m:eq:coupledthree}.

This is a schematic: it fixes the mechanism and the shape of the state, and the
models actually built on it are developed \fwd{rupture}{in the chapter on
compartment rupture}. Two of its features are worth isolating even so, because
they survive into every instance and neither is visible in
\cref{m:eq:coupledgenerator} as written. The continuous variable is
\emph{impulsive}: \cref{m:eq:mediumjump} is a jump map, not a vector field, so
the natural description is a flow punctuated by resets rather than a single
differential equation, and the generator acquires a corresponding jump term in
$y$. And the quantity released, $X_{\tau-}$, is the state of the chain at a
stopping time rather than a parameter of the model: what the rupture finds is
whatever the conditioned process happens to hold, so its law is the
quasi-stationary problem of \cref{m:sec:condmeans} evaluated at a random time.
That is why systems of this kind are approached with the generating-function and
first-step machinery of \cref{m:sec:methods} rather than closed by simulation
alone. \Cref{m:fig:sawtooth} shows one realisation.

\begin{figure}[ht]
\centering
\includegraphics[width=0.86\textwidth]{rupture_sawtooth.pdf}
\caption{One exact event-driven realisation of the schematic
\cref{m:eq:mediumdecay,m:eq:mediumjump}, with $\lambda=0.55$, $\mu=0.25$,
$\rho=0.08$, $X_0=5$, $\delta=0.30$ and $c=0.4$; so that the picture continues
past the first rupture, the compartment is re-seeded with $X_0$ individuals after
each one. Upper panel: the compartment contents $X_t$, growing between ruptures
and collapsing at them. Lower panel: the medium $y(t)$, decaying exponentially
between ruptures and receiving the impulsive release $c\,X_{\tau-}$ at each.
Dotted verticals mark the rupture times, aligned across the panels. The releases
differ in size because each is the compartment's state at the moment its
catastrophe clock happened to fire, and a rate proportional to the contents
biases that state upwards, a fuller compartment rupturing sooner. That is why
\cref{m:eq:mediumjump} cannot be replaced by a fixed increment.}
\label{m:fig:sawtooth}
\end{figure}
